\documentclass[aps,amsmath,amssymb,onecolumn,superscriptaddress,nofootinbib]{revtex4-2}
\usepackage[utf8]{inputenc}

\setcitestyle{super}

\usepackage{graphicx}
\usepackage{dcolumn}
\usepackage{bm}
\usepackage{amsmath}
\usepackage{booktabs}
\usepackage{gensymb}
\usepackage{placeins}
\usepackage[usenames,dvipsnames]{color}
\usepackage[normalem]{ulem}

\newcommand{\comment}[1]{\textcolor{Mahogany}{#1}}
\newcommand{\edit}[1]{\textcolor{OrangeRed}{#1}}
\newcommand{\delete}[1]{\textcolor{OrangeRed}{\sout{#1}}}

% --- Supplementary numbering: Supplementary Note / Fig. S# / Table S# / Eq. S# ---
\renewcommand{\thesection}{Supplementary Note \arabic{section}}
\renewcommand{\thesubsection}{\arabic{section}.\arabic{subsection}}
\renewcommand{\thefigure}{S\arabic{figure}}
\renewcommand{\thetable}{S\arabic{table}}
\renewcommand{\theequation}{S\arabic{equation}}

\begin{document}

\title{Supplementary Information\\[4pt]
\large Two structural states in liquid water revealed by unsupervised classification}
\author{Yiqi Yao}
\author{Diya Sanghi}
\author{Akanksha D. Chokshi}
\author{Minwoo Choi}
\author{Haiqin Wang}
\author{Bilin Zhuang}
\email{bzhuang@hmc.edu}
\affiliation{Department of Chemistry, Harvey Mudd College, 301 Platt Boulevard, Claremont, California 91711}
\affiliation{Yale-NUS College, Singapore 138527, Singapore}
\date{\today}

\maketitle

\noindent
This Supplementary Information provides the extended methods, the full set of
clustering-robustness analyses, and the data and reproducibility notes that
support the main text. Supplementary Notes~1--4 are referenced from the main
text where relevant. Throughout, figures and tables of this document are
prefixed with ``S'' (e.g.\ Fig.~S1, Table~S1); references of the form ``Fig.~1''
without the prefix point to the main-text figures. All quantitative results are
reported for the TIP4P/2005 model at $T=-20\,^{\circ}$C ($253.15$~K) unless
stated otherwise, the condition near which the two populations coexist in
comparable proportions.

\tableofcontents

%=======================================================================
\section{Extended Methods}
\label{sn:methods}

This note expands the Materials and Methods of the main text with the full
simulation protocol and clustering settings, at a level sufficient for
reproduction. The structural order parameters and the per-cluster
structure-factor computation are defined in the main-text Methods and are not
repeated here.

\subsection{Molecular dynamics simulations}
\label{sn:md}

We carry out classical molecular dynamics (MD) simulations of $N=1024$ water
molecules in the canonical (NVT) ensemble with the OpenMM
package. The primary analyses use the rigid,
non-polarizable TIP4P/2005 and
TIP5P models, which reproduce strong tetrahedral
ordering; the polarizable SWM4-NDP model is used as an
additional cross-check (\ref{sn:data}). For each model the
box volume is fixed to reproduce the experimental density at the corresponding
temperature, and simulations are run at $5\,^{\circ}$C increments from
$-30\,^{\circ}$C to $+30\,^{\circ}$C using a Langevin integrator with a
$1$~fs timestep and a friction coefficient of $20$~ps$^{-1}$. After $10$~ns of
equilibration, configurations are sampled every $10$~ps over a $50$~ns
production run. We perform $20$ independent runs per condition, yielding
$N_{\mathrm{samples}} = 20{,}480$ molecular configurations per temperature and
model. The Schottky temperatures, at which the two populations are
approximately equal, are $T_{s=1/2}\approx 237.8$~K for TIP4P/2005 and
$T_{s=1/2}\approx 255.5$~K for TIP5P; $T=-20\,^{\circ}$C
therefore lies close to the TIP4P/2005 coexistence point, which is why we use it
as the reference condition.

\begin{table}[t]
  \centering
  \caption{\textbf{Simulation conditions.} Water models, ensemble, and sampling
  used in this work. Densities are set to the experimental value at each
  temperature; $T_{s=1/2}$ is the Schottky temperature at which the two
  populations are approximately equal.}
  \label{tab:conditions}
  \begin{tabular}{lccccc}
    \toprule
    Model & Type & $T$ range & $T_{s=1/2}$ & Runs & Configs/condition \\
    \midrule
    TIP4P/2005 & rigid, non-pol. & $-30$ to $+30\,^{\circ}$C & $237.8$~K & 20 & $20{,}480$ \\
    TIP5P      & rigid, non-pol. & $-30$ to $+30\,^{\circ}$C & $255.5$~K & 20 & $20{,}480$ \\
    SWM4-NDP   & polarizable     & $-20\,^{\circ}$C          & ---       & --- & cross-check\textsuperscript{a} \\
    \bottomrule
  \end{tabular}
  \\[2pt]
  {\footnotesize \textsuperscript{a}Only a single labelled frame is available for
  SWM4-NDP (\ref{sn:data}).}
\end{table}

\subsection{Clustering algorithms}
\label{sn:clustering}

We benchmark three strategies, all acting on the standardized four-dimensional
feature space.

\emph{K-Means} partitions the data into $k=2$
disjoint clusters by iterative nearest-centroid assignment; centroids are
initialized with K-Means++ from scikit-learn using
$n_{\mathrm{init}}=20$ restarts and a $300$-iteration maximum.

\emph{Gaussian mixture model (GMM)} fits a two-component
mixture by expectation--maximization, providing soft probabilistic assignments;
this choice is motivated by the near-Gaussian order-parameter distributions
predicted by the two-state framework. We additionally fit
$k=2$--$4$ components and compare them using the Bayesian and Akaike information
criteria (BIC, AIC) and the silhouette score (\ref{sn:robustness}).

\emph{Likelihood-ratio classification (main pipeline).} Our main classifier
fits a two-component, full-covariance Gaussian mixture to the standardized
features and assigns each molecule by the class-conditional log-likelihood ratio
$\Lambda(\mathbf{x})=\log p_{\mathrm{LFTS}}(\mathbf{x})-\log p_{\mathrm{DNLS}}(\mathbf{x})$:
LFTS for $\Lambda\geq\tau$, DNLS for $\Lambda\leq-\tau$, and transitional for
$|\Lambda|<\tau$, with $\tau=2$. Prior to fitting, a small fraction of gross
outliers is removed so that the two Gaussian components settle on the two dense
cores rather than being distorted by sparse extreme points. Because the reject band is
symmetric about $\Lambda=0$, it is unbiased toward the minority LFTS state, and
the LFTS/DNLS assignment---being fixed by the sign of $\Lambda$---is independent
of $\tau$.

\emph{Purely density-based clustering (not used).} For comparison we also
applied purely density-based clustering---DBSCAN and HDBSCAN---directly to the
feature space. Because the crossover between the two states is a densely
populated region of overlap rather than a low-density gap, these methods do not
cleanly isolate the two structural states, and we therefore do not use them for
the main classification.

%=======================================================================
\section{Clustering algorithm comparison and robustness}
\label{sn:robustness}

\emph{Per-state order-parameter distributions.}
Figures~\ref{fig:op_kmeans} and~\ref{fig:op_gmm} show the per-state
distributions of the four order parameters, together with the $q$--$\zeta$
scatter, for the K-Means and GMM baselines; the corresponding distributions for
the likelihood-ratio classifier used in the main text are given in main-text
Fig.~1. Both baselines recover the same LFTS/DNLS split: the separation is
sharpest in $\zeta$ but is also clear in $q$ and $S_k$.

\emph{All methods agree, and the likelihood-ratio classifier separates best.}
K-Means and the two-component GMM place the LFTS/DNLS boundary in essentially
the same region of feature space as the likelihood-ratio classifier (main-text
Fig.~1), confirming that the two-state split is a property of the data rather
than of any single algorithm. The silhouette scores are $0.2418$ (K-Means) and
$0.2144$ (GMM) for these every-molecule baselines and $0.34$ for the
likelihood-ratio classifier, which improves separation by first setting aside
the ambiguous transition molecules.

\emph{Hyperparameter selection.}
The main-text likelihood-ratio classifier has a single free parameter, the
reject half-width $\tau$. Because the LFTS/DNLS assignment is set by the sign of
$\Lambda$, $\tau$ controls only the transitional fraction---a broad range of
$\tau$ yields the same two states, with more borderline molecules set aside as
$\tau$ grows---and we use $\tau=2$, which sets aside $\approx 9.8\%$ of molecules.

\emph{Number of clusters.}
Figure~\ref{fig:cluster_number} reports BIC, AIC, and the silhouette score
versus the number of GMM components $k$. BIC and AIC decrease monotonically, with
the largest single drop between $k=1$ and $k=2$, indicating that structural
bimodality is the dominant feature of the data; the absence of a sharp minimum
is expected for thermally broadened, overlapping populations. The silhouette
score resolves the ambiguity, peaking at $k=2$ and declining for $k\geq 3$,
confirming that two components provide the most parsimonious description and that
additional components fragment rather than resolve physical states.

\begin{figure}[t]
  \centering
  \includegraphics[width=\columnwidth]{images/SI/op_kmeans.png}
  \caption{\textbf{K-Means: per-state order-parameter distributions and
  $q$--$\zeta$ scatter} for TIP4P/2005 water at $T=-20\,^{\circ}$C. Probability
  distributions of $q$, LSI, $S_k$, and $\zeta$ for the LFTS (purple) and DNLS
  (orange) states (left), with the classified $q$--$\zeta$ scatter (right).
  K-Means assigns every molecule, so there is no transition-zone (noise)
  population. The bimodal separation is sharpest in $\zeta$ but is also evident
  in $q$ and $S_k$; LSI is in nm$^2$ and $\zeta$ in nm.}
  \label{fig:op_kmeans}
\end{figure}

\begin{figure}[t]
  \centering
  \includegraphics[width=\columnwidth]{images/SI/op_gmm.png}
  \caption{\textbf{GMM: per-state order-parameter distributions and $q$--$\zeta$
  scatter} for TIP4P/2005 water at $T=-20\,^{\circ}$C. As in
  Fig.~\ref{fig:op_kmeans} but for the two-component Gaussian mixture model
  (LFTS, teal; DNLS, brown). The soft probabilistic assignment yields slightly
  more diffuse boundaries than K-Means but recovers the same two-state split.}
  \label{fig:op_gmm}
\end{figure}

% --- DBSCAN/HDBSCAN density-method figures removed from the paper; kept in source for later ---
\iffalse
\begin{figure}[t]
  \centering
  \includegraphics[width=\columnwidth]{images/SI/op_hdbscan.png}
  \caption{\textbf{HDBSCAN$\rightarrow$GMM: per-state order-parameter
  distributions and $q$--$\zeta$ scatter} for TIP4P/2005 water at
  $T=-20\,^{\circ}$C (LFTS, greyish-pink; DNLS, greyish-green; transition-zone
  molecules removed by the HDBSCAN denoising step in grey). Removing the transition zone
  before the GMM fit sharpens the separation, as in the likelihood-ratio
  pipeline of the main text (Fig.~1).}
  \label{fig:op_hdbscan}
\end{figure}

\begin{figure}[t]
  \centering
  \includegraphics[width=0.92\columnwidth]{images/SI/method_comparison.png}
  \caption{\textbf{The two-state separation is robust to the choice of
  clustering algorithm.} Classification of TIP4P/2005 water at
  $T=-20\,^{\circ}$C in the $q$--$\zeta$ plane by K-Means, GMM, and the
  HDBSCAN$\rightarrow$GMM pipeline (LFTS blue, DNLS red). All three place the
  LFTS/DNLS boundary in the same region of feature space; the
  HDBSCAN$\rightarrow$GMM panel additionally shows the transition-zone molecules
  (grey) removed before the GMM fit. The likelihood-ratio result is shown
  in main-text Fig.~1. Silhouette scores: $0.2418$ (K-Means), $0.2144$
  (GMM), $0.2572$ (HDBSCAN$\rightarrow$GMM).}
  \label{fig:method_comparison}
\end{figure}

\begin{figure}[t]
  \centering
  \includegraphics[width=\columnwidth]{images/SI/dbscan_heatmap.png}
  \caption{\textbf{DBSCAN hyperparameter grid search.} \emph{Left}, the
  two-component GMM silhouette after noise removal (colour field) over the
  neighbourhood radius $\varepsilon$ and MinPts, with white contours marking the
  fraction of molecules discarded as transition-zone noise. \emph{Right}, the
  corresponding silhouette--discard tradeoff, coloured by MinPts, with the
  efficiency frontier (the highest silhouette attainable at each discard level).
  Discarding more of the transition zone raises the silhouette at the cost of
  sample size. The operating point (square; $\varepsilon=0.06$,
  $\mathrm{MinPts}=15$; $\approx 23\%$ discarded) defines the DBSCAN
  density-based cross-check.}
  \label{fig:dbscan_heatmap}
\end{figure}

\begin{figure}[t]
  \centering
  \includegraphics[width=\columnwidth]{images/SI/hdbscan_heatmap.png}
  \caption{\textbf{HDBSCAN hyperparameter grid search} (TIP4P/2005,
  $T=-20\,^{\circ}$C). As in Fig.~\ref{fig:dbscan_heatmap}, but for the HDBSCAN
  denoising step over \texttt{min\_cluster\_size} and \texttt{min\_samples} (both
  on log axes): the left panel shows the two-component GMM silhouette after
  removal (colour field) with white contours marking the fraction discarded as
  noise, and the right panel shows the silhouette--discard tradeoff with its
  efficiency frontier. The operating point (square; \texttt{min\_cluster\_size}~$=8$,
  \texttt{min\_samples}~$=8$; $\approx 15\%$ discarded) reproduces the noise mask
  of the HDBSCAN$\rightarrow$GMM classification.}
  \label{fig:hdbscan_heatmap}
\end{figure}
\fi
% --- end removed density-method figures ---

\begin{figure}[t]
  \centering
  \includegraphics[width=0.9\columnwidth]{images/SI/cluster_number.png}
  \caption{\textbf{Model-selection diagnostics identify two clusters.} BIC, AIC,
  and mean silhouette score for GMMs with $k=1$--$4$ components fitted to
  TIP4P/2005 water at $T=-20\,^{\circ}$C. BIC and AIC fall monotonically with the
  largest drop between $k=1$ and $k=2$; the silhouette score peaks at $k=2$ and
  declines for $k\geq 3$. Together the diagnostics select $k=2$ as the most
  parsimonious description.}
  \label{fig:cluster_number}
\end{figure}

\FloatBarrier

%=======================================================================
\section{Robustness across temperature and water model}
\label{sn:generality}

The temperature and model dependence of the two-state signature is presented in
the main text (Fig.~3) and summarized here for completeness. Across the
supercooled-to-ambient range ($-30$ to $+30\,^{\circ}$C) the FSDP of the LFTS
cluster at $k_{T1}$ strengthens monotonically on cooling as the LFTS fraction
grows, while the DNLS peak near $k_{D1}$ remains essentially stable. The same
qualitative behaviour holds for all three water models examined---TIP4P/2005,
TIP5P, and SWM4-NDP---each of which produces a clear FSDP in the LFTS cluster and
a stable ordinary-liquid peak in the DNLS cluster, with peak intensities
ordered TIP5P~$>$~TIP4P/2005~$>$~SWM4-NDP. We caution that the SWM4-NDP
trajectories yield far fewer labelled configurations than the rigid models
(\ref{sn:data}), so its per-cluster $S(k)$ is noisier and
should be read as indicative rather than quantitative.

%=======================================================================
\section{Data, reproducibility, and caveats}
\label{sn:data}

\emph{Sampling and frame counts.}
Cluster labels are computed on a subsample of $20$ frames per condition (the
clustering subsample drawn from the production trajectories), corresponding to
$N=20{,}480$ molecular configurations at each temperature for the rigid models.
Structure-factor uncertainty bands are the per-frame standard error of the mean
over these frames.

\emph{SWM4-NDP.}
The polarizable SWM4-NDP trajectories provide only a single labelled frame at
$T=-20\,^{\circ}$C; its per-cluster $S(k)$ therefore carries no frame-to-frame
error estimate and is presented only as a qualitative cross-check of model
generality.

\emph{Units.}
Distances are taken directly from MDTraj and are expressed in nanometres. In
particular $\zeta$ takes values of order $0$--$0.12$~nm and LSI is in nm$^2$;
earlier internal plots that labelled $\zeta$ in {\AA}ngstr\"om were mislabelled.

\emph{Order-parameter dynamic range.}
The four order parameters differ in how strongly they separate the two states.
The separation is sharpest in $\zeta$ and clear in LSI and $q$, whereas $S_k$
remains within $\sim\!1\%$ of unity for essentially all liquid configurations
(the four nearest-neighbour distances are nearly uniform, so the normalized
nearest-neighbour distance variance that defines $S_k$ is small) and therefore
contributes comparatively little to the cluster assignment. We retain it for completeness and for
comparability with earlier order-parameter studies.

\emph{Cluster identity.}
In the flat per-molecule label files, the LFTS cluster is the one with the
larger mean $\zeta$; in the structure-factor caches the integer cluster
identifiers are reassigned by population and are not stable across temperatures.
We therefore identify the LFTS cluster per condition by its enhanced
FSDP-minus-$D_1$ contrast---the larger value of
$\langle S\rangle_{\mathrm{FSDP}} - \langle S\rangle_{D_1}$ over the reduced-$k$
windows $k\,r_{\mathrm{OO}}/2\pi\in[0.72,0.86]$ and $[0.95,1.10]$. The two
criteria agree on the same physical population (high-$\zeta$, high-$q$,
FSDP-bearing, minority at $-20\,^{\circ}$C with a classified fraction
$\approx 0.36$).

\emph{Code and data availability.}
The order-parameter calculations, clustering pipelines, structure-factor
computation, and the scripts that generate every figure in the main text and
this Supplement are available in the project repository.

\FloatBarrier

\end{document}
